\documentclass[11pt]{article}

\usepackage[utf8]{inputenc}
\usepackage[T1]{fontenc}
\usepackage[protrusion=true,expansion=false]{microtype}
\usepackage[htt]{hyphenat}
\usepackage{amsmath,amssymb}
\usepackage{booktabs}
\usepackage{array}
\usepackage{multirow}
\usepackage{xcolor}
\usepackage{listings}
\usepackage[numbers,sort&compress]{natbib}
\usepackage[colorlinks=true,linkcolor=blue!50!black,citecolor=blue!50!black,urlcolor=blue!50!black]{hyperref}
\usepackage[margin=1in]{geometry}

\newcommand{\ztop}{z_{\mathrm{top}}}
\newcommand{\thp}{\theta}
\newcommand{\softmax}{\operatorname{softmax}}
\newcommand{\code}[1]{\texttt{#1}}
\newcommand{\passk}[1]{pass@#1}

\lstset{
  basicstyle=\ttfamily\small,
  breaklines=true,
  columns=fullflexible,
  keepspaces=true,
  frame=single,
  framesep=4pt,
  xleftmargin=4pt,
  showstringspaces=false,
}

\title{\bf Repetition penalties are not a well-defined intervention:\\
gauge non-invariance and a closed-form corruption threshold\\
in logit-space repetition penalties}

\author{Peter Hollows}
\date{}

\begin{document}
\maketitle

\begin{abstract}
The repetition penalty is a standard decoding knob: one scalar $\thp$, set to around $1.1$--$1.3$
to suppress repetition and assumed to mean the same thing from one model to the next. As
implemented in HuggingFace Transformers (the CTRL divide-if-positive / multiply-if-negative rule
on \emph{raw} logits, inherited in identical form by vLLM and llama.cpp), it is not the clean,
portable knob it is taken to be, for one structural reason. The sign-branch is the community's
own fix for an old bug, in which the naive divide raised the probability of negative-logit
tokens; but the softmax is shift-invariant, so a model's logit zero-point is left unconstrained
by the training objective, and branching on the sign reads exactly that unconstrained
coordinate. The accepted
fix is thus itself ill-posed, with two consequences. Within a single model, the same $\thp$ is
not a single intervention: a behaviorally invisible shift that changes nothing at $\thp=1$ flips
$50$--$94\%$ of greedy tokens at $\thp=1.3$ (five models up to 7B, base and RLHF), and across
models a fixed $\thp$ does not transfer, since real checkpoints sit on different sides of the
sign boundary. Where the penalty acts, it deterministically swaps the model's top token for the
runner-up at a closed-form threshold, $g < \ztop(1-1/\thp)$, which at a routine $\thp=1.3$ drops
schema-valid JSON output to $0\%$. Both effects replicate inside vLLM and llama.cpp through
each stack's own sampler. Both vanish if the penalty is applied to
normalized log-probabilities, an operator HuggingFace already ships but leaves off by default
and applies after the penalty. The fix requires no new machinery; the contribution is the
diagnosis.
\end{abstract}

\section{Introduction}
\label{sec:intro}

Open-ended neural text generation is prone to degeneration: greedy and low-temperature
decoding fall into repetitive loops in which the model re-emits the same token or phrase
indefinitely \citep{holtzman2020curious,welleck2020unlikelihood}. A standard, cheap
countermeasure is the \emph{repetition penalty}: before sampling, the logits of tokens
that have already appeared in the context are pushed down, so that recurrence is
discouraged. Introduced in CTRL \citep{keskar2019ctrl}, the \emph{multiplicative} (sign-branched
divide/multiply) form of the penalty is shipped as \code{repetition\_penalty} in HuggingFace
Transformers \citep{wolf2020transformers} and is carried, in the same form, into vLLM and
llama.cpp's \code{repeat\_penalty}; a separate \emph{subtractive} family
(presence/frequency penalties, exposed by the major hosted APIs and also by vLLM) behaves
differently and, as we note in Section~\ref{sec:related} (Table~\ref{tab:stacks}), is immune
to the first effect we identify. Our results concern the multiplicative CTRL form. Throughout,
our main experiments use the HuggingFace implementation; both effects are additionally
replicated inside vLLM and llama.cpp through each stack's own sampler
(Section~\ref{sec:implications}). Practitioners reach for it reflexively,
set it to a small value
(commonly $1.1$--$1.3$), and reason about it as a clean monotone dial: larger
$\thp$ means ``more anti-repetition,'' the setting is assumed comparable across
models, and the intervention is assumed harmless to everything that is not a degenerate
loop.

This paper shows that, for the penalty as it is actually implemented, all three of those
assumptions fail in measurable ways. A \code{repetition\_penalty} value
tuned on one model is not the same operation on another, so it does not transfer in any
principled sense; and at a routine $\thp=1.3$ the operator can deterministically corrupt
structured output, dropping a model's rate of valid, schema-conformant JSON to zero. Both
effects, and the fix for both, trace to a single structural decision, itself introduced as a
fix. CTRL's original penalty divided every seen logit by $\thp$, but dividing a
\emph{negative} logit by $\thp>1$ moves it toward zero and \emph{raises} the token's
probability, so the penalty backfires on exactly the tokens it should suppress; this was
reported to HuggingFace in 2019 \citep{hf2302}. The remedy, now standard, is to branch on the
\emph{sign} of each raw logit (multiply negatives, divide positives), so a seen token's
probability always falls. That remedy works at the level it was aimed at. Our observation is
one level up:
branching on a raw logit's sign reads a quantity the model is free to choose arbitrarily,
because the softmax is invariant to a global additive shift of the logits. The accepted
fix for the naive penalty is itself ill-posed. We make this precise and draw out two
consequences.

\paragraph{A1: the penalty is not gauge-invariant (Section~\ref{sec:a1}).} The simplest
form of the result lives \emph{inside a single model}. Adding a constant $c$ to
every logit is a provable no-op at $\thp=1$ (we confirm it is token-for-token invisible under
both greedy and seeded sampling); yet with the penalty on, two such provably-equivalent gauges
of the \emph{same} model disagree on $50$--$94\%$ of greedy tokens at $\thp=1.3$. The two gauges
differ only by an additive constant that has no effect on the output distribution. The cross-model reading
follows as a corollary: because a model's logit zero-point is never trained or constrained, real
checkpoints sit at different points on the sign boundary, so a fixed $\thp$ is a different
operation on each. (Multiplicative penalties on un-normalized logits would fail to transfer
across models for mundane reasons too, since logit scales and calibration differ; the
within-model result is the one that isolates the gauge, rather than scale, as the cause.)

\paragraph{A2: where it acts, it is a closed-form corruption threshold
(Section~\ref{sec:a2}).} Structured output (JSON, source code, quoted strings, list
markers) contains \emph{grammar-obligatory} repetition: the same delimiters
($\code{\}}$, $\code{]}$, $\code{,}$, $\code{"}$, indentation) must recur, and the penalty
cannot tell ``the grammar requires this token'' from ``the model is stuck in a loop.'' We
derive the exact condition under which the penalty flips a confident, correct token to the
runner-up---$g < \ztop(1-1/\thp)$---and show it predicts the realized flip with
balanced accuracy $\ge 0.999$ on two independently trained 7B code models, with every flip
landing on the pre-penalty second choice.

We deliberately separate what is verified from what is not. The original framing of A2
predicted that corruption would strike the model's \emph{most confident} tokens first and
would be sharply \emph{code-specific}; both are refuted by our data
(Section~\ref{sec:a2-refuted}). Reporting these negative sub-results is part of the
contribution: the surviving statement (a gap-driven threshold that corrupts all token types
similarly but lands disproportionately on the delimiters that structured output cannot avoid) is
more precise and better supported than the original.

Both results concern the same operator, and both are removed by the same one-line change:
apply the penalty to normalized log-probabilities rather than raw logits
(Section~\ref{sec:fix}). Normalization is itself shift-invariant, so it erases the gauge
$c$; and it makes every value non-positive, so the sign-branch stops being a branch at
all. We confirm this directly: under the normalized operator the A1 flip rate collapses to
$0.000$ at every $\thp$ on all three models, and the A2 structural corruption rate drops by
roughly $27\times$. HuggingFace already ships this normalization operator but leaves it
disabled by default and, even when enabled, runs it \emph{after} the penalty, too late to
help. The remediation is a reorder of existing machinery; the observation is that it is not the
default.

\section{The operator under test}
\label{sec:operator}

We study the operator exactly as implemented in HuggingFace Transformers
(\code{Repetition\-Penalty\-Logits\-Processor}), which is the de facto reference and matches the
CTRL formulation \citep{keskar2019ctrl}. Let $z \in \mathbb{R}^{V}$ be the logits the model
emits for the next token over a vocabulary of size $V$, and let $S \subseteq \{1,\dots,V\}$
be the set of token ids that have already appeared in the context (prompt plus generated
text). The penalty with strength $\thp \ge 1$ transforms only the seen tokens, by
branching on the sign of each logit:
\begin{equation}
\label{eq:operator}
z'_i =
\begin{cases}
z_i / \thp & i \in S \text{ and } z_i \ge 0,\\[2pt]
z_i \cdot \thp & i \in S \text{ and } z_i < 0,\\[2pt]
z_i & i \notin S.
\end{cases}
\end{equation}
The next token is then drawn from $\softmax(z')$ (or chosen as $\arg\max z'$ under greedy
decoding). The verbatim implementation is:
\begin{lstlisting}[language=Python]
score = torch.gather(scores, 1, input_ids)
# if score < 0 multiply by penalty (more negative); else divide (smaller positive)
score = torch.where(score < 0, score * self.penalty, score / self.penalty)
scores_processed = scores.scatter(1, input_ids, score)
\end{lstlisting}

Two implementation facts matter for everything that follows. First, the operator acts on
\emph{raw} logits: there is no normalization before it. Second, HuggingFace does ship a
normalization processor,
\begin{lstlisting}[language=Python]
class LogitNormalization(LogitsProcessor):
    def __call__(self, input_ids, scores):
        return scores.log_softmax(dim=-1)
\end{lstlisting}
but \code{GenerationConfig.renormalize\_logits} defaults to off, and when it is enabled the
processor is appended to the \emph{end} of the processor list, after the repetition
penalty. So in the default configuration, and even in the non-default ``renormalize''
configuration, the penalty sees un-normalized logits. This is precisely the degree of
freedom we exploit in Section~\ref{sec:a1} and remove in Section~\ref{sec:fix}.

The sign-branch is not incidental: it is the community's fix for the naive divide (which raised
the probability of negative-logit tokens; Section~\ref{sec:intro}, \citealp{hf2302}). What
matters here is that the two branches meet at $z_i=0$ with \emph{different slopes}, so the
operator's behavior turns on which side of $z_i=0$ each seen logit lands, and $z_i=0$ is a point
the softmax lets the model place anywhere (Section~\ref{sec:a1}).

\section{A1: gauge non-invariance}
\label{sec:a1}

\subsection{Theory}
\label{sec:a1-theory}

\paragraph{The gauge.} For any constant $c \in \mathbb{R}$ and the all-ones vector
$\mathbf{1}$,
\begin{equation}
\label{eq:shiftinv}
\softmax(z + c\,\mathbf{1})_i
= \frac{e^{z_i + c}}{\sum_j e^{z_j + c}}
= \frac{e^{c} e^{z_i}}{e^{c}\sum_j e^{z_j}}
= \softmax(z)_i .
\end{equation}
The output distribution is invariant to a global additive shift of the logits. Adding $c$
to every logit is identical to adding $c$ to every component of the final linear layer's
bias ($\code{lm\_head.bias} \mathrel{+}= c$); the greedy choice is therefore unchanged
exactly, and sampled output is unchanged token-for-token under a fixed random
seed. The shift $c$ is a pure \emph{gauge}: $c$ parameterizes an orbit of logit
vectors that the softmax collapses to one and the same distribution, so any well-defined
function of the model's \emph{behavior} must be invariant to it. We use ``gauge'' in exactly
this sense (an unobservable redundancy a well-posed decoding operator ought to ignore), and
nothing in the argument relies on the term beyond it. At $\thp=1$ the penalty is the identity
and this invisibility is exact.

\paragraph{The sign-branch breaks it.} The penalty in Eq.~\eqref{eq:operator} is not a
function of the distribution; it is a function of the raw logits, and specifically of
their signs. Apply the gauge first. A previously-seen token with logit $x$ now has logit
$x + c$, and the penalty subtracts from it
\begin{equation}
\label{eq:drop}
\Delta(x;c) =
\begin{cases}
(x+c)\left(1 - \tfrac{1}{\thp}\right) & x + c \ge 0 \quad(\text{divide branch}),\\[4pt]
|x+c|\,(\thp - 1) & x + c < 0 \quad(\text{multiply branch}).
\end{cases}
\end{equation}
Every \emph{competing} (unseen) token also shifts by $c$, so the constant $c$ cancels in
all the comparisons that decide the next token, \emph{except} through the nonlinear,
sign-dependent transform applied to the penalized tokens. The drop $\Delta$ is not
shift-invariant: it depends on where $x+c$ sits relative to the kink at zero, which the
gauge $c$ controls.

Concretely, consider a high-probability repeated token. With $c = +5$ it sits deep in the
positive/divide regime, where $\Delta$ grows with $x+c$, and it is heavily penalized. With
$c = -5$ the same token can be pushed just below zero into the multiply branch, where
$|x+c|$ is small and $\Delta \approx 0$: the penalty nearly vanishes on the
token it was meant to suppress. Same model, same $\thp$, same prompt; two gauges that
are provably indistinguishable at $\thp=1$ now implement two different interventions.

\subsection{Design}
\label{sec:a1-design}

We implement the gauge as a logits transform $z \mathrel{+}= c$ applied immediately before
the penalty in a hand-written decode loop. This is mathematically identical to shifting
the \code{lm\_head} bias, but makes the ordering explicit and sidesteps GPT-2's tied,
bias-free output head. The penalty uses the exact HuggingFace semantics of
Eq.~\eqref{eq:operator} over all previously-seen ids (prompt and generation). We run three
independently trained models (\code{gpt2}, \code{gpt2-large} \citep{radford2019gpt2}, and
\code{pythia-2.8b} \citep{biderman2023pythia}) in fp32, over a grid of
$\thp \in \{1.0, 1.02, 1.05, 1.1, 1.15, 1.3\}$ and $c \in \{-5,-3,-1,0,1,3,5\}$, with
16 fixed prompts and 200 generated tokens each.

\paragraph{No-op control.} The validity gate is that the gauge must be invisible when the
penalty is off. At $\thp=1$ we require the generated sequence to be token-for-token
identical across all $c$, for every prompt, under both greedy and seeded sampling. Greedy
invariance is trivial; the seeded-sampling leg (a shared RNG stream across gauges) is the
stronger check, since it probes the full sampled distribution rather than only its argmax.
We do not enforce bitwise equality in the harness; we \emph{verify} softmax shift-invariance
empirically, and it holds: this control passed exactly, $32/32$ (prompt, mode) pairs
identical, with $0$ of $3200$ greedy tokens differing between $c=+5$ and $c=-5$ at
$\thp=1$. Because the gauge is added inside the same decode loop that runs every $\thp>1$
condition, dtype, RNG, and code path are differenced out within each comparison; any
divergence at $\thp>1$ is therefore caused by the gauge $\times$ sign-branch interaction, not by
numerics or a bug.

\paragraph{Endpoint.} The quantity that directly operationalizes the claim ``a provable
no-op changes generation'' is the per-step argmax \emph{flip rate}: the fraction of greedy
decoding steps at which the $c=+5$ and $c=-5$ runs (both provable no-ops at $\thp=1$)
choose different tokens. Unlike a repetition ratio, the flip rate cannot saturate: it is
$0$ exactly when behavior is unchanged and rises as the gauge bites. We report it as a
function of $\thp$, with a per-prompt bootstrap 95\% confidence interval ($10^4$
resamples) at $\thp=1.3$.

\subsection{Results}
\label{sec:a1-results}

The flip rates are in Table~\ref{tab:a1}. At $\thp=1$ every model flips exactly
$0\%$ of tokens: the no-op control passes exactly. At $\thp=1.3$ the \emph{same model}, run under two
provably-equivalent gauges ($c=+5$ and $c=-5$), differs on $50\%$ (gpt2), $91\%$
(pythia-2.8b), and $94\%$ (gpt2-large) of greedy tokens, bootstrap intervals far above zero
throughout; the model itself is unchanged. The effect is not confined to small base models: it replicates at 7B and is
undiminished by post-training, with the no-op gate exactly $0$ and flip rates of $0.871$ on
Qwen2.5-7B and $0.897$ on its instruction-tuned/RLHF variant Qwen2.5-7B-Instruct
(Table~\ref{tab:a1}), if anything stronger than on the small models. The synthetic shift
$c=\pm5$ used here is a controlled sweep (it maps the response surface and pins the no-op
gate), and the constant $5$ is arbitrary. Below we use a gauge that involves no
arbitrary constant: each model's \emph{own median} logit, a re-centring it could equally
have shipped, which flips more tokens (Table~\ref{tab:zeropoint}). The flip rate
rises from zero as soon as the penalty engages ($0.61$--$0.78$ at $\thp=1.02$); its
magnitude is non-monotone in $\thp$ for the small loop-prone gpt2, and the effect manifests as
changed content rather than changed repetition rate (audit in App.~\ref{app:audit}).

\begin{table}[t]
\centering
\caption{A1: argmax flip rate between the $c=+5$ and $c=-5$ gauges (both provable no-ops at
$\thp=1$). The gate column confirms the no-op is exact; the $\thp=1.3$ column shows
the same provably-equivalent shift selects different greedy output once the penalty is on.
The effect replicates at 7B and is undiminished by RLHF (last two rows).}
\label{tab:a1}
\begin{tabular}{lccc}
\toprule
Model & flip rate ($\thp{=}1$) & flip rate ($\thp{=}1.3$) & 95\% CI \\
\midrule
gpt2-large            & $0.000$ & $0.941$ & $[0.916,\,0.959]$ \\
pythia-2.8b           & $0.000$ & $0.905$ & $[0.876,\,0.933]$ \\
gpt2                  & $0.000$ & $0.497$ & $[0.333,\,0.655]$ \\
\midrule
Qwen2.5-7B (base)     & $0.000$ & $0.871$ & $[0.832,\,0.906]$ \\
Qwen2.5-7B-Instruct   & $0.000$ & $0.897$ & $[0.864,\,0.927]$ \\
\bottomrule
\end{tabular}
\end{table}


\paragraph{Real checkpoints differ in zero-point.} Nothing about the gauge $c$
requires a synthetic shift. A model's logit zero-point is never an optimization target and is
not constrained by any loss term; it is determined incidentally by initialization, the (often
absent or untied) output bias, weight decay, and training dynamics, and so it differs across
checkpoints and architectures. To quantify this, we
measured where each of our five real checkpoints sits relative to the sign boundary
the penalty branches on, at decode time (Table~\ref{tab:zeropoint}). The spread is wide:
the fraction of already-seen tokens whose logit is positive (i.e.\ which the penalty
\emph{divides} rather than \emph{multiplies}) ranges from $0.09$ on gpt2 to $0.99$ on
Qwen2.5-Coder-7B. At one and the same \code{repetition\_penalty=1.3}, then, the operator
divides the logits of $\sim$$9\%$ of seen tokens on gpt2 (so it mostly \emph{multiplies}) and
$\sim$$99\%$ on Qwen2.5-Coder (so it almost always \emph{divides}): on these two checkpoints
the sign-branch sends the penalty down opposite paths, with no synthetic shift
applied to either. gpt2 is the extreme case (a small model with a tied, bias-free head and a
median top-1 logit of $-97.5$, far below the others), but the effect is not an artifact of that
one outlier: even setting gpt2 aside, the divide-fraction still spans $0.75$ to $0.99$ across
the four remaining models, a wide enough range that a fixed $\thp$ applies a materially
different penalty to each.

\begin{table}[t]
\centering
\caption{Where five real checkpoints sit relative to the penalty's sign boundary, measured at
decode time ($c=0$, $\thp=1$). The fraction of seen-token logits that are positive (the
divide branch) spans $0.09$--$0.99$ ($0.75$--$0.99$ excluding the small gpt2 outlier), so a
fixed $\thp$ is a materially different operation on each. The last column is the per-model A1
flip rate between the model's raw gauge and its own median-centred gauge (a gauge it could
equally have shipped) at the \emph{natural} offset magnitude.}
\label{tab:zeropoint}
\begin{tabular}{lccc}
\toprule
Model & frac.\ seen logit $>0$ (divide) & median top-1 logit & flip @ natural $\thp{=}1.3$ \\
\midrule
gpt2              & $0.091$ & $-97.5$ & $0.569$ \\
gpt2-large        & $0.746$ & $11.9$  & $0.781$ \\
starcoder2-7b     & $0.828$ & $15.5$  & $0.858$ \\
pythia-2.8b       & $0.938$ & $17.6$  & $0.815$ \\
Qwen2.5-Coder-7b  & $0.986$ & $24.2$  & $0.796$ \\
\midrule
\multicolumn{3}{l}{\emph{pooled across the five checkpoints}} & $0.764$ \\
\bottomrule
\end{tabular}
\end{table}

The natural offsets are large. Re-centring each model's
logits by its own median (a behaviorally invisible gauge it could equally have been trained
into) and comparing the resulting greedy trajectory to the raw one (both provable no-ops at
$\thp=1$) flips $57$--$86\%$ of tokens per model at $\thp=1.3$ ($76.4\%$ pooled across the
five), \emph{more} than a synthetic $c=5$ shift ($31.6\%$ pooled). The gauge non-invariance is therefore not an artifact
of an adversarially large injected shift; it is at least as strong at the offset magnitudes
real checkpoints already exhibit. Consequently \code{repetition\_penalty=1.3}
already denotes different interventions on different real models, and a setting tuned on
one does not transfer to another in any principled sense, because the two models do not
share the coordinate system the penalty operates in.

\section{A2: the corruption threshold}
\label{sec:a2}

\subsection{The closed form}
\label{sec:a2-theory}

Structured output contains repetition that the grammar \emph{requires}: a JSON object must
close every \code{\{} with a \code{\}}, separate items with \code{,}, and quote keys with
\code{"}; Python must re-emit indentation and colons. These delimiters have, by
construction, appeared before, so they are in the penalized set $S$, and they are typically
the model's confident top choice at the position where the grammar demands them. The
penalty does not distinguish grammar-obligatory recurrence from degenerate looping, so it
penalizes them.

Consider such a position. Let the correct token be the pre-penalty top choice with logit
$\ztop > 0$ (so it is in the divide branch), and let the runner-up be a gap $g = \ztop -
z_2 \ge 0$ below it. Assume the runner-up is unpenalized (it is not a delimiter that just
occurred). The penalty maps $\ztop \mapsto \ztop/\thp$ but leaves the runner-up at
$z_2 = \ztop - g$. The emitted (greedy) token flips from the correct top choice to the
runner-up exactly when
\begin{equation}
\label{eq:threshold}
\frac{\ztop}{\thp} < \ztop - g
\quad\Longleftrightarrow\quad
g < \ztop\!\left(1 - \frac{1}{\thp}\right).
\end{equation}
The right-hand side is the \emph{vulnerable gap}: any correct token whose lead over the
runner-up is smaller than $\ztop(1-1/\thp)$ is deterministically corrupted. The
threshold is exact (not probabilistic) under greedy decoding, and it is a hard threshold in
$\thp$: the flip occurs discontinuously, not gradually.

\subsection{Design}
\label{sec:a2-design}

We test Eq.~\eqref{eq:threshold} on \code{bigcode/starcoder2-7b}
\citep{lozhkov2024starcoder2}, a model whose domain makes grammar-obligatory delimiters
abundant. Decoding is greedy with $\thp \in \{1.0,1.1,1.2,1.3,1.4,1.5\}$ over 18 fixed
prompts (6 JSON, 6 Python, and 6 prose as a control with no grammar-obligatory repetition),
256 generated tokens each. The model runs in bf16, but at every step we upcast the final
logits to fp32 before the penalty, the argmax, and the threshold evaluation, so that the
formula test is not blurred by bf16 rounding. At each position we record the pre-penalty
top-1 id and $\ztop$, the runner-up id and $z_2$, the gap $g$, indicator flags
(\,\code{top1\_seen}, \code{top1\_positive}, \code{top2\_seen}, \code{top1\_structural}\,),
and the post-penalty emitted id, from which we derive \code{flip} and
\code{emitted\_is\_runnerup}. A token is tagged \emph{structural} if its decoded form is
made entirely of delimiter characters $\{\}[]()",:;=$ or whitespace/indentation.

\paragraph{Validity gate and a numerical subtlety.} At $\thp=1.0$ the penalty is the
identity, so the flip rate must be exactly $0$ at every position; this passes. Reaching it
required one fix worth flagging for replicators: a naive bf16 \code{topk}-versus-\code{argmax}
comparison reports spurious ``flips'' at ties, because the two can break an exact tie
differently. The harness therefore defines the pre-penalty top-1 by the same \code{argmax}
used post-penalty; with that, the $\thp=1.0$ gate is clean.

\subsection{What is verified}
\label{sec:a2-verified}

We evaluate Eq.~\eqref{eq:threshold} on the \emph{clean} positions where it is exact: the
top-1 token is seen and positive (divide branch) and the runner-up is unseen (unpenalized).
Pooling over $\thp \in \{1.1,\dots,1.5\}$ on the code prompts gives $n=3973$ clean
positions ($25.9\%$ of all penalized code positions; the remaining $74.1\%$ are the
two-sided case where the runner-up is itself penalized, treated below), of which $2225$
actually flip.

On this slice the threshold is, by derivation, an algebraic identity for greedy argmax: the
penalty maps $\ztop\to\ztop/\thp$ and leaves the runner-up untouched, so $g<\ztop(1-1/\thp)$
\emph{is} the flip condition. We therefore do not read the predicate accuracy
(balanced accuracy $0.999$; TPR $1.000$, TNR $0.999$; the only departures from $1.000$ are
bf16 rounding) as the empirical content. The empirical content is threefold and is \emph{not}
guaranteed by the algebra:
\begin{itemize}
\item $100\%$ of the $2225$ flips emit the pre-penalty runner-up. When the penalty
corrupts a token, it does not scatter probability diffusely; it deterministically promotes
the specific second-place token. A diffuse-corruption alternative would not produce this.
\item The $\thp=1$ gate holds (0 flips), so every flip is penalty-induced, not an
artifact.
\item The threshold is crossed at scale: a penalty of
$\thp\in[1.1,1.5]$ flips $\sim$$31\%$ of all greedy code tokens (not only clean ones), so
real structured decoding lands in this regime routinely.
\end{itemize}
Where the penalty acts, it
either leaves the choice alone or deterministically replaces the model's top token with its
runner-up, at a point computable in closed form from the pre-penalty logits. The two-sided
positions ($74.1\%$) obey a more general two-sided inequality (both top-1 and runner-up
penalized) that we state in Section~\ref{sec:limitations} but do not separately tabulate; the
exact closed form above is the one-sided regime.

\paragraph{Replication on a second code model.} To check that the threshold is a property of
the operator rather than of one checkpoint's logit geometry, we re-ran the identical harness
on \code{Qwen2.5-Coder-7B} \citep{hui2024qwen25coder}, an independently trained code model.
The result is, if anything, sharper: on $n=3584$ clean penalized positions ($2484$ flips)
the closed form achieves balanced accuracy $1.000$ (TPR $1.000$, TNR $1.000$),
$100\%$ of flips emit the pre-penalty runner-up, and the $\thp=1.0$ validity gate
passes. The same gap-driven signature recurs (Section~\ref{sec:a2-refuted}): flipped tokens
have lower mean $\ztop$ ($23.9$) than non-flipped ($28.8$). The corruption threshold is thus
demonstrated on two independent 7B code models.

\paragraph{What is refuted.}\label{sec:a2-refuted} We pre-registered two stronger claims about the threshold, and our
own data refute both. It does \emph{not} corrupt the most confident tokens first: corruption is
\emph{gap-driven}, so flipped tokens have \emph{lower} mean $\ztop$ than non-flipped ones (e.g.
$17.5$ vs $22.1$ on StarCoder2, and the same ordering on Qwen2.5-Coder). And it is not sharply
code-specific: the code-vs-prose structural-flip ratio is real but model-dependent ($2.3\times$
on StarCoder2, $5.6\times$ on Qwen2.5-Coder), and prose is corrupted at a comparable overall
per-token rate. The penalty corrupts all token types at a similar rate and lands disproportionately on
delimiters only because structured output exposes more of them. The full audit, with CIs, is in
Appendix~\ref{app:audit}. What survives is the narrower statement: the penalty
deterministically flips the greedy token to the runner-up exactly when $g < \ztop(1-1/\thp)$: a
gap-driven threshold that corrupts all token types similarly and so lands disproportionately on
the grammar-obligatory delimiters structured output cannot avoid.

\section{The fix: normalize before penalizing}
\label{sec:fix}

Both effects trace to one decision (branching on the sign of an un-normalized logit), and
both are removed by one change: apply the penalty to normalized log-probabilities,
$\ell = \log\softmax(z)$, instead of to raw logits $z$. This is less ``remove a degree of
freedom'' than ``use the right coordinate.'' Among all the gauges $z+c\mathbf{1}$ that the
softmax identifies, $\ell$ is the canonical representative, the one in which the model's own
distribution is actually written ($\sum_i e^{\ell_i}=1$); the operator should read there rather
than at a zero-point the training objective leaves unconstrained. Two properties of $\ell$ do the work.

\paragraph{(a) Normalization erases the gauge.} $\log\softmax$ is itself shift-invariant:
$\log\softmax(z + c\mathbf{1}) = \log\softmax(z)$, since the additive $c$ cancels between
the logit and the log-partition. Penalizing $\ell$ therefore cannot depend on $c$, and the
A1 divergence is eliminated by construction: the operator becomes invariant to the
coordinate choice it currently reads.

\paragraph{(b) The penalty becomes tempering, with a positive identity.} Every log-probability
satisfies $\ell_i \le 0$, so under Eq.~\eqref{eq:operator} every penalized token takes the
\emph{multiply} branch, $\ell_i \mapsto \thp\,\ell_i$, moving each log-probability uniformly
further from zero. There is no kink whose location depends on an arbitrary zero-point, and no
token can sit on the ``cheap'' side of a sign boundary the way the A1 mechanism exploits. And
the operation now has a name: $\thp\,\ell_i = \thp\log p_i = \log p_i^{\thp}$, so in probability
space it is $p_i \mapsto p_i^{\thp}$ (before renormalization): the fixed operator
tempers each seen token's probability toward zero by a common power. This is a
monotone, well-defined
suppression measured in log-probability units, a coordinate-free quantity, and hence at least
\emph{dimensionally} comparable across models (we do not claim a fixed $\thp$ is numerically
optimal across models, only that it denotes the same kind of operation). The normalized penalty does act on
log-probabilities ($O(1)$ in magnitude) rather than raw logits ($O(10)$), so the same $\thp$ is
a gentler push, which is exactly why A2 corruption collapses below.

\paragraph{Remediation.} The remediation is a design principle, not a patch to one library:
a decoding-time penalty should act on a shift-invariant quantity (normalized
log-probabilities, or any transform from which the gauge has been removed), and the
normalization must run before the penalty, not after. For the operator at hand, that means
penalizing $\log\softmax(z)$ rather than $z$. It applies unchanged to every stack carrying the
multiplicative CTRL form (HuggingFace, vLLM, llama.cpp; Table~\ref{tab:stacks}), since they
share the operator, and it is nearly free: no new machinery, only a reorder. HuggingFace makes
the point concrete: it already ships exactly this normalization as \code{LogitNormalization}
(\code{scores.log\_softmax(dim=-1)}), but off by default (\code{renormalize\_logits=None}) and,
when enabled, appended \emph{after} the penalty, so it normalizes scores the penalty has already
modified.\footnote{The related idea of normalizing \emph{after} the logits processors is known in a
different setting: because processors can leave beam scores un-normalized, HuggingFace
re-normalizes them (via \code{log\_softmax}) so beam-search comparison is fair \citep{hf10638}.
Applying normalization \emph{before} the penalty, and to the repetition penalty specifically, is
the point here.} The
fix is thus reachable in stock machinery yet absent from the default path, which is the finding.
Its one cost is that, being gentler per unit $\thp$, the normalized penalty does not inherit
existing tuned values: $\thp$ must be re-set; in exchange, the same $\thp$ denotes the same
operation on every model.

\paragraph{Why fix it rather than abandon it?} One could instead switch to a shift-invariant
subtractive penalty (Table~\ref{tab:stacks}) or a purpose-built alternative such as the LZ
penalty \citep{ginart2025lz}; those are complementary. We focus on the in-place, one-line
correction because the multiplicative CTRL form is the shipped default and existing tuned
presets are written against it.

\paragraph{Demonstration.} We re-ran both experiments with a single change, inserting
$z \leftarrow \log\softmax(z)$ immediately before the penalty (exposed as a \code{--fix}
flag) and holding everything else identical, and the prediction holds (Table~\ref{tab:fix}). For A1,
the gauge-induced flip rate, which the un-normalized operator drives to $0.50$--$0.94$ at
$\thp=1.3$, collapses to $0.000$ at every $\thp$ on all three models: the gauge $c$
becomes invisible everywhere, not just at $\thp=1$, exactly as shift-invariance requires.
For A2, the code structural-flip rate falls from $0.138$ to $0.005$ ($\sim$$27\times$
lower), and the closed-form threshold of Eq.~\eqref{eq:threshold} stops being predictive
(balanced accuracy $0.999 \to 0.794$), the intended outcome: there is no longer a sign-branch
threshold for it to predict.\footnote{The residual $0.794$ is asymmetric (TPR $1.000$, TNR
$0.587$): under the normalized operator the multiply branch sends nearly every penalized
log-probability across the predicate's right-hand side, so $g<\ztop(1-1/\thp)$ (still evaluated
on the raw-logit $\ztop$ and $g$, and no longer the condition governing the flip) fires almost
everywhere while actual flips have collapsed. The formula stops \emph{discriminating} flips
because there are almost none left to discriminate; it is not a competing threshold.}

\begin{table}[t]
\centering
\caption{The fix (apply the penalty to $\log\softmax(z)$). A1 gauge divergence is eliminated
on all three models; A2 delimiter corruption is reduced $\sim$$27\times$ and the threshold
formula ceases to fire.}
\label{tab:fix}
\begin{tabular}{llcc}
\toprule
& metric & raw logits (default) & normalized (fix) \\
\midrule
\multirow{3}{*}{A1: flip rate at $\thp{=}1.3$}
 & gpt2-large  & $0.941$ & $0.000$ \\
 & pythia-2.8b & $0.905$ & $0.000$ \\
 & gpt2        & $0.497$ & $0.000$ \\
\midrule
\multirow{2}{*}{A2: StarCoder2-7B}
 & code structural-flip rate & $0.138$ & $0.005$ \\
 & threshold balanced accuracy & $0.999$ & $0.794$ \\
\bottomrule
\end{tabular}
\end{table}

The reason corruption nearly vanishes (rather than the penalty being switched
off outright) is instructive. After normalization a confident correct delimiter has
log-probability close to $0$, so the multiply branch ($\ell_i \mapsto \ell_i\cdot\thp$)
barely moves it; the penalty's effect is now proportional to how \emph{improbable} a token
already is, so it pushes down overused tokens while leaving confident,
grammar-required ones essentially untouched. This is the intended behavior of a repetition
penalty; the un-normalized sign-branch does not provide it.

\paragraph{The corruption matters end-to-end, and the fix recovers it.} The results so far are
mechanistic (per-token flips, threshold accuracy). To confirm the per-token corruption changes
something a user observes, we measured two end-to-end quality metrics on \emph{complete}
generations from a single code model (Qwen2.5-Coder-7B), at the common settings
$\thp\in\{1.0,1.1,1.3\}$, under the raw operator versus the fix (Table~\ref{tab:downstream}).
On the JSON metric: prompting for a complete JSON object against one
of six fixed schemas ($n=48$ prompts) and validating with \code{json.loads} plus a type check,
the raw penalty's valid-rate falls from $1.00$ at $\thp=1$ to $0.00$ at
$\thp=1.3$: at a routine \code{repetition\_penalty=1.3} the operator yields no
parseable, schema-valid JSON, while under the fix it stays at $1.00$ across all $\thp$. A
HumanEval \passk{1} run shows the same shape (raw $0.079\to0.000$; fix flat at ${\sim}0.079$);
we report only the \emph{relative} trend there, because our lightweight completion-extraction
harness depresses the absolute baseline far below the model's true HumanEval, so the absolute
numbers are a harness artifact and not a model fact. This is one model and a small probe, but
it connects the per-token threshold (A2) to a user-facing failure. The exact
$\thp$ at which a given model's structured output collapses will itself vary with the model's
zero-point (Table~\ref{tab:zeropoint}); the mechanism is what generalizes, and the fix removes
it.

\begin{table}[t]
\centering
\caption{End-to-end quality under the raw vs.\ normalized penalty (single model,
Qwen2.5-Coder-7B, greedy, complete generations; JSON $n=48$ over six schemas). At a routine
$\thp=1.3$ the raw operator yields no valid JSON; the fix preserves validity. For HumanEval
only the relative $\thp$-trend is meaningful (the absolute baseline is depressed by our
extraction harness; see text), so we report it as corroborating, with JSON conformance as the
clean headline. This is an existence demonstration of end-to-end impact, not a universal rate.}
\label{tab:downstream}
\begin{tabular}{llccc}
\toprule
Metric & operator & $\thp{=}1.0$ & $\thp{=}1.1$ & $\thp{=}1.3$ \\
\midrule
\multirow{2}{*}{JSON schema valid-rate}
 & raw & $1.00$ & $1.00$ & $0.00$ \\
 & fix & $1.00$ & $1.00$ & $1.00$ \\
\midrule
\multirow{2}{*}{HumanEval \passk{1} (trend only)}
 & raw & $0.079$ & $0.043$ & $0.000$ \\
 & fix & $0.079$ & $0.085$ & $0.079$ \\
\bottomrule
\end{tabular}
\end{table}

\paragraph{The fix is a usable penalty, not a disabled one.} A natural worry is that
normalize-before-penalize removes the corruption only by weakening the penalty into a no-op.
It does not: on the loop-prone gpt2 and pythia-2.8b, the normalized penalty still breaks
degenerate loops, with the repetition rate falling monotonically in $\thp$ (gpt2
$0.83\to0.63$, pythia-2.8b $0.80\to0.60$ over $\thp\in[1.0,1.5]$). It is, as expected,
\emph{gentler per unit $\thp$} than the raw operator, which reaches a repetition rate of
${\sim}0.01$--$0.1$ by $\thp=1.3$: because the normalized penalty acts on bounded
log-probabilities rather than gauge-dependent raw logits, the same $\thp$ is a smaller push,
so the operative range shifts upward. This both confirms the fix retains anti-degeneration
efficacy and quantifies the re-tuning caveat above.

\section{Implications}
\label{sec:implications}

\paragraph{Both effects replicate inside vLLM and llama.cpp.} The results above are measured on
the HuggingFace implementation; the same operator ships in vLLM and llama.cpp
(Table~\ref{tab:stacks}), and we re-ran A1 and A2 inside each, through each stack's own sampler
and its own knob: vLLM 0.8.5.post1 (the operator in \code{vllm/main} at commit \code{e196268} is
mathematically identical) and llama.cpp at master commit \code{4fc4ec55}. The A1 gauge probe
passes its validity gate in both stacks (exactly $0$ of $3{,}200$ greedy tokens flip at
$\thp=1$) and reproduces the divergence at $\thp=1.3$ on gpt2-large: flip rate $0.939$ in vLLM
and $0.941$ in llama.cpp, against $0.941$ on HuggingFace (Table~\ref{tab:a1}). The A2 JSON task
(Qwen2.5-Coder-7B, same schemas and validator as Table~\ref{tab:downstream}) goes from $100\%$
valid at $\thp\in\{1.0,1.1\}$ to $16.7\%$ at $\thp=1.3$ in both stacks (HF: $0\%$; the residual
valid outputs trace to bf16 numerics). llama.cpp's default 64-token penalty window does not
change the outcome: greedy JSON completes inside the window, and its output is byte-identical
to whole-context penalization on all 144 paired generations.\footnote{Measuring inside
llama.cpp also surfaced an unrelated defect: \code{--repeat-last-n -1} is documented as
``context size'' but the sampler clamps the value to $0$, silently disabling the penalty
(\code{src/llama-sampler.cpp}, \code{std::max(penalty\_last\_n, 0)}).} Scripts and raw outputs
are in \code{code/smoke\_*} and \code{results/smoke\_*} alongside the main harness.

\paragraph{Penalty settings are not portable.} Because the operator reads an unconstrained
gauge (A1), a \code{repetition\_penalty} value carries no fixed meaning across models or
checkpoints. Default cookbooks, leaderboard configurations, and tuned presets that pin a
single value implicitly assume a comparability the operator does not provide. Under the
normalized operator of Section~\ref{sec:fix}, the penalty's magnitude is expressed in
log-probability units and is at least dimensionally comparable across models, a
precondition for any value to transfer.

\paragraph{Structured-output corruption is predictable.} Practitioner
folklore that repetition penalties ``break JSON'' or ``mangle code'' (reported anecdotally
around vLLM and guided-decoding tooling, and a motivation for grammar-constrained decoding
\citep{willard2023outlines}) has, in A2, an exact account: the penalty flips a confident,
correct delimiter to its runner-up whenever $g < \ztop(1-1/\thp)$, and this reaches
end-to-end output: at a routine \code{repetition\_penalty=1.3} it takes schema-valid JSON
output to $0\%$ (Table~\ref{tab:downstream}). Practitioners who must keep a repetition penalty
on for structured generation can either exempt grammar-obligatory tokens from the penalized
set, or adopt the normalized operator, under which the sign-branch corruption
disappears and JSON validity is restored to $100\%$. Related reports are current: 2026 issues against
Gemma~4, in both the reference stack \citep{gemma622} and vLLM \citep{vllm40080}, document JSON
generation collapsing into repetition that grammar-constrained decoding does not prevent (the
repeated words are valid string content) and that \code{repeat\_penalty} at $1.0$, $1.15$, and
$1.5$ does not fix. That failure mode is distinct from ours (there the penalty
fails to help; here it causes the failure); we note the fielded mitigation is to detect and
truncate the loop
\citep{vllm40099}. Either way, the
multiplicative penalty is not a reliable instrument for structured generation.

\paragraph{Operational footprint.} The two effects have different
operational footprints, and the difference favors A1. A2's deterministic corruption applies
wherever an unconstrained, low-temperature multiplicative penalty meets structured output:
grammar-constrained decoding \citep{willard2023outlines} neutralizes it (the masked,
grammar-invalid runner-up cannot be emitted), and raising temperature dilutes the deterministic
flip into a merely probabilistic one. A1, by contrast, is \emph{not} rescued by either:
the gauge shifts the entire logit vector \emph{before} temperature scaling and independently of
any token mask, so it perturbs the whole distribution at any temperature and under any
constraint set. The non-portability of penalty values (A1) therefore reaches every deployment
of the multiplicative penalty, while the structured-output corruption (A2) is concentrated in the
common but bounded regime of unconstrained greedy/low-temperature structured generation.

\paragraph{The general form.} Any decoding-time operator that branches on, or thresholds
against, an absolute (non-shift-invariant) function of the raw logit inherits a dependence on a
coordinate the training objective leaves unconstrained. The CTRL sign-branch is the most widely
deployed instance;
the dividing line is absolute-zero-point dependence, not ``reading magnitude'' per se. Operators
that read only shift-invariant quantities are immune by construction: top-$k$, top-$p$/nucleus,
min-$p$ \citep{nguyen2024minp}, typical sampling \citep{meister2023typical}, and contrastive
search \citep{su2023contrastive} all act on the softmax distribution or on logit
\emph{differences}, where the gauge cancels.

\section{Related work}
\label{sec:related}

\paragraph{Anti-degeneration decoding.} Repetition penalties sit alongside truncated sampling
\citep{fan2018hierarchical,holtzman2020curious}, unlikelihood training
\citep{welleck2020unlikelihood}, contrastive decoding \citep{su2022contrastive}, and penalty
redesigns such as the information-theoretic LZ penalty \citep{ginart2025lz} as tools against
degeneration. That literature asks \emph{how much} to penalize, or proposes a better penalty; we
ask, prior to that, whether the standard operator is well-defined, and show a fixed $\thp$ is
not a fixed intervention.

\paragraph{Multiplicative vs.\ subtractive penalties (scope of A1).} Two penalty families are
deployed in practice, and the gauge dependence distinguishes them (Table~\ref{tab:stacks}).
The \emph{multiplicative} CTRL form (Eq.~\eqref{eq:operator}) sign-branches on the raw logit and
is the form shipped by HuggingFace and by llama.cpp's \code{repeat\_penalty} \citep{llamacpp};
it is the one we study, and it is subject to A1. The \emph{subtractive} presence/frequency
penalties exposed by the major hosted APIs \citep{openai_api_penalties} (and also by vLLM
\citep{kwon2023vllm}) instead subtract a constant (and, for the frequency variant, a
count-scaled term) from the logit of seen tokens: $z_i \mapsto z_i - \alpha$. A
constant subtraction commutes with the gauge ($\,(z+c\mathbf{1})_i - \alpha = (z_i-\alpha)+c$,
and the $c$ again cancels in softmax), so the subtractive family is \emph{shift-invariant and
immune to A1}; the count-scaled frequency term is likewise gauge-invariant since the counts do
not depend on $c$. A2 has no analogue for the subtractive form at all, because there is no
sign-branch. This is why we scope the title's claim to the multiplicative operator: the
ubiquity is real, but it is the ubiquity of \emph{that} operator, not of every penalty knob.

\begin{table}[t]
\centering
\caption{Penalty form by stack, and exposure to the two effects. The multiplicative CTRL
form (sign-branch on raw logits) is subject to A1 gauge non-invariance and the A2 threshold;
the subtractive presence/frequency form subtracts a constant and is shift-invariant, hence
immune to both. All three multiplicative stacks are measured directly (HuggingFace in the main
experiments; vLLM and llama.cpp in the cross-stack replication, Section~\ref{sec:implications}).}
\label{tab:stacks}
\begin{tabular}{llccc}
\toprule
Stack / knob & Penalty form & Normalizes first? & A1? & A2? \\
\midrule
HF \code{repetition\_penalty}        & multiplicative (CTRL) & no (default) & yes & yes \\
llama.cpp \code{repeat\_penalty}     & multiplicative (CTRL) & no & yes & yes \\
vLLM \code{repetition\_penalty}      & multiplicative (CTRL) & no & yes & yes \\
Hosted-API \code{presence\_penalty}  & subtractive (constant) & n/a & no & no \\
\code{frequency\_penalty} (API/vLLM) & subtractive (count-scaled) & n/a & no & no \\
\bottomrule
\end{tabular}
\end{table}

\paragraph{Logit geometry.} Work on glitch and under-trained tokens \citep{land2024magikarp}
exploits a model's embedding geometry through specific anomalous tokens; our gauge dependence is
instead a \emph{global} additive freedom in the output logits, independent of any particular
token.

\paragraph{Constrained decoding.} Grammar-constrained decoding \citep{willard2023outlines} routes
around structured-output corruption by masking invalid tokens at the sampling layer; we instead
characterize it in closed form and fix the operator itself. The two specific contributions (the
gauge non-invariance of the sign-branched penalty and the closed-form delimiter threshold) are,
to our knowledge, new.

\section{Limitations}
\label{sec:limitations}

We study the specific HuggingFace/CTRL multiplicative operator of Eq.~\eqref{eq:operator}; the
subtractive presence/frequency penalties (Section~\ref{sec:related}, Table~\ref{tab:stacks})
are shift-invariant and escape A1, and stacks that already normalize before penalizing are
likewise exempt. The A1 flip-rate sweep spans five models up to 7B (three small base models
plus Qwen2.5-7B and its instruction-tuned/RLHF variant, on which the effect is undiminished;
Table~\ref{tab:a1}), and the zero-point realism measurement (Table~\ref{tab:zeropoint})
covers all five checkpoints; we have not pushed beyond 7B or to the largest frontier models. The
cross-stack replication (Section~\ref{sec:implications}) uses vLLM 0.8.5.post1 (with the
operator verified unchanged in \code{vllm/main} at \code{e196268}) and a single pinned
llama.cpp master commit; its A1 leg covers gpt2-large only. A2
is demonstrated on two independently trained 7B code models (StarCoder2-7B and
Qwen2.5-Coder-7B), which share the closed-form threshold and the gap-driven signature but
differ in the magnitude of code-vs-prose concentration ($2.3\times$ versus $5.6\times$), so
that secondary quantity should be read as model-dependent. Our A2 analysis isolates the
``clean'' positions where Eq.~\eqref{eq:threshold} is exact (top-1 seen and positive,
runner-up unseen); positions where both the top choice and the runner-up are penalized obey a
two-sided version of the inequality that we do not separately tabulate here. Finally, our
core endpoints are mechanistic (per-token flips, threshold accuracy); we connect them to
end-to-end quality on two tasks (Table~\ref{tab:downstream}), but a broader downstream
evaluation (more models, more schemas, a canonical HumanEval harness for clean absolute
\passk{1}, and prose-quality metrics) is left to future work.

\section{Conclusion}
\label{sec:conclusion}

The repetition penalty is treated as a clean, portable dial that leaves everything but repetition
alone. As actually implemented (a sign-branched divide/multiply on raw logits), it is none of these. It is not
gauge-invariant: a provably invisible shift of the logit zero-point, which already differs across
real models, turns the same nominal $\thp$ into a different intervention,
flipping half to nearly all greedy tokens at $\thp=1.3$. And where it acts it is a
closed-form corruption threshold, $g < \ztop(1-1/\thp)$, that deterministically replaces a
confident correct token with its runner-up, verified at balanced accuracy $0.999$, with
every flip landing on the runner-up. Both effects share one root cause and one remedy:
normalize before penalizing. The operator that does this already ships in the library; it is
off by default and runs after the penalty. The penalty is not the well-defined knob it is
assumed to be, and making it one requires only a reorder of existing machinery.

\section*{Reproducibility}
All experiments are inference-only and run on a single 24--48\,GB GPU. The small A1 models run
in fp32; the 7B models use bf16 (\code{run\_a1.py --dtype bfloat16}) to fit the card, but the
gauge shift and argmax decision are computed in fp32, so the no-op gate is exact regardless of
model dtype. The pinned environment
(\code{torch 2.6.0+cu124}, \code{transformers 5.11.0}), the public model revisions, and the
exact scripts that produce every number and table are included with this paper:
\code{run\_a1.py}/\code{analyze\_a1b.py} (A1, with \code{diag\_divergence.py} for the
divergence diagnostic) and \code{run\_a2.py}/\code{analyze\_a2.py} (A2, with
\code{explore\_delim.py} for the brackets-only re-slice). The fix demonstration
(Table~\ref{tab:fix}) is reproduced by \code{run\_a1.py --fix} / \code{run\_a2.py --fix}
($\to$ \code{results/a1b\_fix/}, \code{results/a2\_fix/}), and the second-model replication
by \code{run\_a2.py} on Qwen2.5-Coder-7B ($\to$ \code{results/a2\_qwen/}). The end-to-end
quality probe (Table~\ref{tab:downstream}) is produced by
\code{run\_a2\_downstream.py}/\code{analyze\_a2\_downstream.py}, the fix loop-suppression check
by \code{run\_a1\_loopcheck.py}, and the zero-point realism measurement
(Table~\ref{tab:zeropoint}) by \code{run\_a1\_zeropoint.py}. The cross-stack replication is
produced by \code{code/smoke\_vllm/} and \code{code/smoke\_llamacpp/} (each with its own README,
pinned environment, and build steps; raw outputs and per-stack reports in
\code{results/smoke\_*}). Pre-registered hypotheses and
frozen decision rules, and a record of which sub-claims were verified, partial, or
refuted, accompany the code.

\appendix

\section{Pre-registered hypotheses: verified, partial, and refuted}
\label{app:audit}

Each experiment was pre-registered with a frozen decision rule before its confirmatory run; the
rules and their raw verdicts ship with the code. The body reports the surviving claims; here we
record the sub-hypotheses that were only partially confirmed or refuted, so the audit trail is
complete without slowing the main argument.

\subsection{A1: measurement channel and the shape of the curve}

\paragraph{The effect is content divergence, not a repetition-rate change.} The original
hypothesis predicted that a weaker-gauge penalty would leave more repetition. At $\thp=1.3$ that
channel is invisible: the penalty breaks loops under \emph{both} gauges, and a distinct-2
endpoint saturates at $1.000$ for every $c$. The divergence instead manifests as different
\emph{content}, which is what the flip-rate endpoint measures and why we adopt it. In a
non-saturated regime (the smallest $\thp=1.02$) the repetition-rate channel is recoverable but
weak: for gpt2-large and pythia-2.8b the repetition rate is ordered in the predicted direction
across $c$ (e.g.\ gpt2-large $0.823 \to 0.771$ from $c=-5$ to $c=+5$), while for gpt2 it is flat
and noisy. The strong, unambiguous signal is content divergence.

\paragraph{Divergence versus $\thp$: off-zero leap, model-size-dependent shape, and a
monotonicity clause that overreached.} The flip rate rises from exactly $0$ at $\thp=1$ to a
large value as soon as the penalty turns on ($0.61$--$0.78$ at $\thp=1.02$, near-ties
tipped by the sign-branch), with per-$\thp$ bootstrap $95\%$ CIs far above zero at every
$\thp>1$ (percentile bootstrap over the $16$ prompts, $10^4$ resamples). The later shape varies
with model size: the larger models climb monotonically toward $\sim$$0.9$, while small gpt2
peaks near $\thp=1.05$ and trends down (as strong loop-breaking forces both gauges into similar
non-repetitive text). Our pre-registered A1 rule included a clause requiring the flip rate to be
\emph{monotone non-decreasing} in $\thp$; that fails for gpt2 and pythia-2.8b, so the frozen rule
returned ``not confirmed'' for those two. We leave the frozen verdict as printed, but the clause was
an overreach: nothing in the claim (``a gauge no-op changes behavior'') entails
that the \emph{amount} of divergence grow monotonically with $\thp$. The substantive clauses
(flip rate jumps off zero; bootstrap CI excludes zero) hold on all three models. We also report
the gpt2 non-monotonicity as directional only: at $n=16$ the bump and its $\thp=1.3$ value have
overlapping CIs.

\subsection{A2: two refuted embellishments and an uninformative control}

\paragraph{``Corrupts the most confident tokens first'': refuted.} The reasoning was that the
vulnerable gap $\ztop(1-1/\thp)$ grows with $\ztop$, so high-confidence tokens should be hit
first. It is incomplete: confident tokens also tend to have \emph{larger} gaps $g$, and the gap
term dominates. Flipped tokens have \emph{lower} mean $\ztop$ than non-flipped tokens: $17.5$
versus $22.1$ over all penalized positions (difference $95\%$ CI $[-4.7,-4.4]$), and $19.9$
versus $23.9$ restricted to delimiters (CI $[-4.4,-3.6]$); the same ordering holds on
Qwen2.5-Coder ($23.9$ versus $28.8$, CI $[-5.1,-4.8]$). Corruption is \emph{gap-driven}: it lands
on tokens the model is \emph{less} sure about. The threshold formula stands; the
prediction about \emph{which} tokens it selects was wrong.

\paragraph{Code-specificity is real but model-dependent.} The structural-token flip rate per
generated token is higher in code than prose, but the ratio varies with the model: on StarCoder2
it is $0.138$ versus $0.060$, only $2.3\times$ ($2.4\times$ on a brackets-only re-slice), below
the pre-registered $\ge 5\times$; on Qwen2.5-Coder it is $0.196$ versus $0.035$, a $5.6\times$
that clears the bar. Two things are stable. First, prose is corrupted at a comparable
\emph{overall} rate (on StarCoder2 prose flips \emph{more} tokens than code, $36\%$
versus $31\%$), so the penalty corrupts all token types at a similar rate and is not syntax-specific.
Second, code concentrates \emph{which} tokens are hit: $19\%$ of code's flips land on delimiters
versus $7\%$ in prose. The concentration on \emph{structural} tokens is because structured output
exposes more grammar-obligatory delimiters to a gap-driven threshold.

\paragraph{The validity ``cliff'' was uninformative as measured.} We had intended to show a sharp
drop in parse validity as $\thp$ crosses the threshold, but StarCoder2's free continuations are
truncated fragments, so \code{json.loads} sits at the floor ($0.00$) even at $\thp=1.0$ and
bracket-balance error versus $\thp$ is noise. We make no cliff claim from it; it was
pre-registered as descriptive and non-gating. (The clean end-to-end measurement on complete
documents is in Section~\ref{sec:fix}, Table~\ref{tab:downstream}.)

\bibliographystyle{plainnat}
\bibliography{refs}

\end{document}
